% sec-02-methods.tex - Section 2, Methods.  DRAFT stage.
% One page maximum of prose.  Three figure slots, Figures 3, 4 and 5, and two
% tables, Tables 6 and 7.
\section{Methods}
\label{sec:methods}

The method is one master prompt, decomposed by the model into sub-prompts that
the model then executes in order, with every generated file committed and pushed
before the next is written. Earlier 2026 papers in this repository reached a
similar effect through a sequence of user prompts, one per stage; the difference
is that the author had to be present at each boundary. Under the master prompt
schedule the author is present once, at the start, and thereafter watches the
branch.

\draftinstr{Stage 7: expand to one full page. Sources are
\appfile{funding/capitalization-plan/prompts/prompt-capital.md},
\appfile{funding/capitalization-plan/sub-prompts/README.md}, and
\appfile{new-trial-system/prompts/prompt-new-trial.md}. State the count of works
built each way: multi-prompt through mid 2026, master prompt from
pdac-funding-applications onward. Give the exact stage count, 8, and the commit
rule, one file one commit.}

\figslot{plantuml-type / activity with fork and join}{9.4}{Figure 3. The eight-stage build as one activity: five diagram stages
fork and\\join on a shared specification, then three paper stages run strictly in
order.}

Figure 3 separates the part of the build that is genuinely concurrent from the
part that is not. Five diagram stages run against one fixed figure plan and must
all close before any LaTeX is written; the three paper stages that follow are
serial, because each reads the whole of its predecessor. Table 6 counts how many
times each method element has been used across separate deposited projects,
which is the basis on which the method is offered as reliable rather than novel.

\draftinstr{Stage 7: build Table 6 from
\appfile{new-trial-system/abstracts/README.md} and the prompts and sub-prompts
directories of \appfile{funding/capitalization-plan} and
\appfile{funding/pdac-funding-applications}. Columns: Method element, Projects
using it, First use, Evidence. Rows: master prompt, AI-generated sub-prompt
schedule, five-platform diagram specification set, draft to full to final
staging, one-file-one-commit real-time push, machine-readable figure reuse,
AI peer review round, DOI deposit per artifact.}

\begin{apptable}
\begin{tabularx}{\textwidth}{>{\raggedright\arraybackslash}p{4.6cm} >{\raggedright\arraybackslash}p{1.7cm} >{\raggedright\arraybackslash}p{2.4cm} >{\raggedright\arraybackslash}X}
\headrow \thead{Method element} & \thead{Projects} & \thead{First use} & \thead{Evidence}\\
\midrule
Master prompt, single turn & 3 & Jul 2026 & \appfile{funding/pdac-funding-applications/prompts} \\
AI-generated sub-prompt schedule & 3 & Jul 2026 & \appfile{funding/capitalization-plan/sub-prompts} \\
Five-platform diagram specification set & 3 & Jul 2026 & Five directories per project \\
Draft, full, final staging & 8 & Mar 2026 & Every paper directory since March \\
One file, one commit, pushed live & 3 & Jul 2026 & Branch history of each build \\
Machine-readable figure reuse & 3 & Jul 2026 & Specifications re-read, not re-derived \\
AI peer review round & 4 & Nov 2025 & \S7 \\
DOI deposit per artifact & 30 or more & 2024 & \appfile{new-trial-system/abstracts/README.md} \\
\end{tabularx}
\end{apptable}
\vspace{-0.6cm}
\tabcap{Table 6. How many separate deposited projects have used each element\\
of the method, and where the evidence for each of those counts is held.}

\figslot{mermaid-type / sequenceDiagram}{8.8}{Figure 4. One generation turn from master prompt to pushed commit,
with the\\four participants, the two return paths, and the point the author can
stop it.}

The second half of the method is storage. Every figure in this paper exists
first as a specification: a perspective statement, a two-line caption, valid
source in the platform's own syntax, a construction table of absolute
coordinates, an edge-routing paragraph, and a value table naming the repository
file each number came from. Figure 4 shows how one such file reaches the branch,
and Figure 5 shows what it holds once it is there.

\figslot{graphviz-type / dot record}{7.0}{Figure 5. Two records for one figure: the machine-readable
specification\\and the raster image, with the six fields only the first returns as
input.}

\draftinstr{Stage 7: build Table 7 from
\appfile{new-trial-system/graphviz/fig-05-figure-storage-record.md}. Columns:
Property, Specification, Raster, Consequence for the next paper. Add the
measured total size of the twenty-five specification files in this build.}

\begin{apptable}
\begin{tabularx}{\textwidth}{>{\raggedright\arraybackslash}p{3.3cm} >{\raggedright\arraybackslash}p{3.5cm} >{\raggedright\arraybackslash}p{3.0cm} >{\raggedright\arraybackslash}X}
\headrow \thead{Property} & \thead{Specification} & \thead{Raster} & \thead{Consequence for the next paper}\\
\midrule
Size on disk & 6 to 9 kilobytes & 180 to 900 kilobytes & The whole figure set fits in one file's space \\
Caption & Two lines, exact & None & Captions are diffable across versions \\
Coordinates & Absolute, stated & None & A node can be added without a re-layout \\
Routing & Stated per edge & None & Overlap claims are auditable \\
Values & Sourced per cell & None & No number in a figure is unattributable \\
Reuse & Read back as input & Re-drawn by hand & The next paper starts from this one \\
\end{tabularx}
\end{apptable}
\vspace{-0.6cm}
\tabcap{Table 7. What a machine-readable figure specification stores that a\\raster
does not, and what each stored field buys the next paper.}

\draftinstr{Stage 8: confirm one page of prose, and that Figures 3, 4 and 5 and
Tables 6 and 7 are each referenced once by name.}
